\documentclass[a4paper,11pt]{article}
        \usepackage[top=15mm,bottom=18mm,left=16mm,right=16mm]{geometry}
        \usepackage{fontspec}
        \usepackage{xeCJK}
        \setmainfont{Times New Roman}
        \setCJKmainfont{Yu Gothic}
        \setmonofont{Consolas}
        \setCJKmonofont{Yu Gothic}
        \usepackage{booktabs}
        \usepackage{longtable}
        \usepackage{tabularx}
        \usepackage{array}
        \usepackage{enumitem}
        \usepackage{hyperref}
        \usepackage{listings}
        \usepackage{xcolor}
        \hypersetup{
          colorlinks=true,
          linkcolor=blue,
          urlcolor=blue,
          pdftitle={live ROS2 launch 入口},
          pdfauthor={Codex}
        }
        \lstset{
          basicstyle=\ttfamily\small,
          breaklines=true,
          columns=fullflexible,
          keepspaces=true,
          frame=single,
          framerule=0.2pt,
          rulecolor=\color{black},
          showstringspaces=false
        }
        \setlength{\parskip}{0.42em}
        \setlength{\parindent}{1em}
        \renewcommand{\arraystretch}{1.15}
        \title{04. live ROS2 launch 入口}
        \author{solar\_ws0129-main}
        \date{2026-07-29}
        \begin{document}
        \maketitle

        \section*{対象}
        \begin{tabularx}{\linewidth}{>{\raggedright\arraybackslash}p{0.18\linewidth} X}
        \toprule
        項目 & 内容 \\
        \midrule
        ファイル & \path|launch/solar_race_live.launch.py| \\
        種別 & ROS 2 launch Python \\
        区分 & launch \\
        役割 & WiFi 文字列 bridge を使わない live 運用の基本 launch。 \\
        起動文脈 & live モード起動時に ros2 launch される。 \\
        \bottomrule
        \end{tabularx}

        \section*{このファイルを読む理由}
        WiFi 文字列 bridge を使わない live 運用の基本 launch。

        \begin{itemize}[leftmargin=1.6em]
        \item profile を読み build\_live\_nodes に委譲する。
\item ノード構成の本体は live\_launch.py にある。
        \end{itemize}

        \section*{呼び出し関係}
        \subsection*{どこから呼ばれるか}
        \begin{itemize}[leftmargin=1.6em]
        \item \path|scripts/solar_control.sh|
\item \path|SolarSim.ps1|
        \end{itemize}

        \subsection*{次に読むと理解が進むファイル}
        \begin{itemize}[leftmargin=1.6em]
        \item \path|mpc_solarcar/live_launch.py|
\item \path|mpc_solarcar/mpc_node.py|
        \end{itemize}

        \subsection*{同一リポジトリ内の主要依存}
        \begin{itemize}[leftmargin=1.6em]
        \item \path|mpc_solarcar/live_launch.py|
\item \path|mpc_solarcar/solar_profile.py|
        \end{itemize}

        \section*{ソース冒頭の把握}
        モジュール docstring または先頭意図の要約:

        モジュール docstring は明示されていない。

        \subsection*{代表コード断片}
        \begin{lstlisting}[language=Python]
from mpc_solarcar.solar_profile import load_profile


def _setup(context):
    profile_yaml = LaunchConfiguration("profile_yaml").perform(context)
    profile_path, cfg = load_profile(profile_yaml)
    return build_live_nodes(profile_path, cfg, use_wifi=False)


def generate_launch_description():
    return LaunchDescription(
        [
            DeclareLaunchArgument("profile_yaml", default_value="config/solar/bwsc_2027_demo.yaml"),
            OpaqueFunction(function=_setup),
        ]
    )
\end{lstlisting}


        \section*{主要構造の読み取り}
        主要関数は generate\_launch\_description。

        \subsection*{主要クラス・関数}
            \begin{longtable}{p{0.07\linewidth} p{0.16\linewidth} p{0.25\linewidth} p{0.40\linewidth}}
            \toprule
            行 & 種別 & 名前 & 補足 \\
            \midrule
            10 & function & \path|_setup| & args: context \\
16 & function & \path|generate_launch_description| &  \\
            \bottomrule
            \end{longtable}

        \subsection*{ファイルを上から読んだときの定義順}
            \begin{longtable}{p{0.07\linewidth} p{0.84\linewidth}}
            \toprule
            行 & 内容 \\
            \midrule
            10 & 関数 \_setup を定義する。 \\
16 & 関数 generate\_launch\_description を定義する。 \\
            \bottomrule
            \end{longtable}

        \subsection*{import 群の詳細}
            \begin{longtable}{p{0.07\linewidth} p{0.33\linewidth} p{0.50\linewidth}}
            \toprule
            行 & import 文 & このファイルで読む理由 \\
            \midrule
            1 & \path|from __future__ import annotations| & 型ヒントの遅延評価を有効にし、自己参照型や forward reference を安全に扱うため。 このファイル内での使用位置は少ないか、間接利用である。 \\
3 & \path|from launch import LaunchDescription| & launch が実行すべき action 群をまとめるため。 このファイル内での主な使用位置は L17。 \\
4 & \path|from launch.actions import DeclareLaunchArgument, OpaqueFunction| & DeclareLaunchArgument や OpaqueFunction など launch action を使うため。 このファイル内での主な使用位置は L19, L20。 \\
5 & \path|from launch.substitutions import LaunchConfiguration| & launch 引数の実行時値を参照するため。 このファイル内での主な使用位置は L11。 \\
6 & \path|from mpc_solarcar.live_launch import build_live_nodes| & live系 node 構成ビルダ から build\_live\_nodes を読み込み、このファイルの内部処理を分担させるため。 実体ファイルは mpc\_solarcar/live\_launch.py。 このファイル内での主な使用位置は L13。 \\
7 & \path|from mpc_solarcar.solar_profile import load_profile| & profile YAML 読込と検証 から load\_profile を読み込み、このファイルの内部処理を分担させるため。 実体ファイルは mpc\_solarcar/solar\_profile.py。 このファイル内での主な使用位置は L12。 \\
            \bottomrule
            \end{longtable}

        \section*{関数・クラスを上から順に解説}
        \subsection*{L10 関数 \path|_setup|}
            \begin{itemize}[leftmargin=1.6em]
              \item 定義: \path|_setup(context)|
              \item docstring: 明示されていない。
              \item このブロックが直接呼ぶ主な関数・メソッド: LaunchConfiguration, build\_live\_nodes, load\_profile, perform
              \item 戻り値の要点: build\_live\_nodes(profile\_path, cfg, use\_wifi=False)
            \end{itemize}
            \begin{enumerate}[leftmargin=1.8em]
            \item profile\_yaml に LaunchConfiguration('profile\_yaml').perform(context) の結果を代入する。
\item (profile\_path, cfg) に load\_profile(profile\_yaml) の結果を代入する。
\item build\_live\_nodes(profile\_path, cfg, use\_wifi=False) を返す。
            \end{enumerate}

\subsection*{L16 関数 \path|generate_launch_description|}
\begin{itemize}[leftmargin=1.6em]
  \item 定義: \path|generate_launch_description()|
  \item docstring: 明示されていない。
  \item このブロックが直接呼ぶ主な関数・メソッド: DeclareLaunchArgument, LaunchDescription, OpaqueFunction
  \item 戻り値の要点: LaunchDescription([DeclareLaunchArgument('profile\_yaml', default\_value='config/solar/bwsc\_2027\_demo.yaml'), OpaqueFunction(function=\_setup)])
\end{itemize}
\begin{enumerate}[leftmargin=1.8em]
\item LaunchDescription([DeclareLaunchArgument('profile\_yaml', default\_value='config/solar/bwsc\_2027\_demo.yaml'), OpaqueFunction(function=\_setup)]) を返す。
\end{enumerate}











        \section*{処理の流れ}
        \begin{enumerate}[leftmargin=1.7em]
        \item launch 引数を受け取り、profile を読み込む。
        \end{enumerate}

        \section*{このファイルの位置づけ}
        この資料群では、\path|launch/solar_race_live.launch.py| を
        \textbf{launch}の中核として扱う。
        したがって、ここで扱う処理は単独で完結せず、
        前段の profile / launch / telemetry と後段の planner / logger / report へ繋がる。

        \end{document}
